\section{A nonlinear activation at every layer, uniformly close to identity}
\label{sec:near}

This model has the appended constant coordinate $(1,\xi)$ and uses the
same activation $\phi$ in all $L-1$ hidden layers. Set $d=2k+1$,
$p=d+1$, $m=L-1$, and $Q_*=Q_M$ from \eqref{eq:Qt}.

\begin{theorem}[All widths for near-identity activations]\label{thm:near}
Fix $k\ge3$ and $L\ge2$. Choose
\begin{equation}\label{eq:epsilon}
 0<\epsilon\le
       \frac1{128m p^{3/2}Q_*^{3/2}},\qquad \alpha=1-\epsilon.
\end{equation}
Suppose $\phi(0)=0$, $\phi$ is $1$-Lipschitz, and
\begin{equation}\label{eq:close-id}
 |\phi(s)-s|\le\epsilon|s|\qquad(s\in\R).
\end{equation}
For the target and fixed input order of Section~\ref{sec:selector}, at
every fixed learner width $N\ge1$, the correct-label posterior
probability on each of the $T_k$ designated rounds tends to zero as
$\sigma\downarrow0$ under \eqref{eq:prior}. Consequently, at every
sufficiently small positive scale, posterior majority makes all $T_k$
mistakes, the strict $2/3$ confidence criterion fails on all $T_k$ rounds,
and posterior sampling makes all $T_k$ mistakes with probability at
least $1-\eta$, for any prescribed $\eta>0$.

At every depth $L\ge k+2$, the target has $C_L(f)=2$ for this activation.
Thus the lower bound applies to every $N\ge C_L(f)$ with no upper width
bound. The activation condition is independent of $N$; the positive
scale threshold may depend on $N,k,L,\phi,\eta$.
\end{theorem}

Two examples are
\[
 \phi(s)=\begin{cases}s,&s\ge0,\\(1-\epsilon)s,&s<0,\end{cases}
 \qquad\text{and}\qquad
 \phi(s)=(1-\epsilon)s+\epsilon\tanh s.
\]
Both are nonlinear for $\epsilon>0$. The second is smooth,
odd, and has derivative between $1-\epsilon$ and one. Our sufficient upper bound on $\epsilon$ in \eqref{eq:epsilon} is of
order $8^{-k}/[(L-1)k^{3/2}]$.
The theorem therefore does not give a growing-target lower bound for
one fixed standard nonlinear activation.

\begin{lemma}[Linearization on a quadratic-energy ball]
\label{lem:linearize}
Replace every hidden activation in $\theta$ by the identity, obtaining
an affine output $A_\theta\cdot(1,\xi)$. Under the theorem's hypotheses,
\begin{equation}\label{eq:linear-error}
 \max_\xi|Z_\theta(\xi)-A_\theta\cdot(1,\xi)|
 \le\epsilon m\sqrt{p/d}\left(\frac{\kappa_{d,N}(\theta)}L\right)^{L/2}.
\end{equation}
Also
\begin{equation}\label{eq:linear-energy}
 \kappa_{d,N}(\theta)\ge
       L\bigl(d\norm{A_\theta}_2^2\bigr)^{1/L}.
\end{equation}
\end{lemma}
\begin{proof}
Since $\phi$ is $1$-Lipschitz and fixes zero, it does not increase
Euclidean vector norm. Telescope through the $m$ activation locations,
replacing one at a time. The replacement error is at most $\epsilon$
times the incoming vector norm, and every remaining map has Lipschitz
constant at most its matrix Frobenius norm. Each summand is therefore
bounded by $\epsilon\sqrt p$ times the product of all $L$ matrix norms,
including the readout. The coefficients in the energy are
$d/N,1,\ldots,1,N$, whose product is $d$. Weighted
arithmetic--geometric mean gives
\[
 \prod_{j=1}^L\norm{M_j}_F
 \le d^{-1/2}\left(\frac{\kappa_{d,N}(\theta)}L\right)^{L/2},
 \qquad (M_1,\ldots,M_L)=(W^1,\ldots,W^{L-1},v).
\]
This proves \eqref{eq:linear-error}. Submultiplicativity also bounds
$\norm{A_\theta}_2$ by this product, proving \eqref{eq:linear-energy}.
\end{proof}

\begin{proof}[Proof of Theorem~\ref{thm:near}]
Fix a history through $q_t$, $k\le t<M$. First construct a cheap
consistent predictor with the wrong sign at $q_{t+1}$. Embed a scalar
path in one neuron of each hidden layer, with positive downstream
weights, whose identity-activation product vector is $P A_t$, where
\[
 P=\frac1{2\alpha^m}.
\]
Condition \eqref{eq:close-id} implies that $\phi$ preserves every
nonzero sign and $|\phi(s)|\ge\alpha|s|$. The resulting nonlinear output
has sign $A_t$ and magnitude at least
$\alpha^m P|A_t(\xi)|\ge1/2$ on the full cube. Its training signs are
correct and its next-query sign is positive.

Balance the $L$ factors against their Gaussian energy coefficients.
Explicitly, let $c_1=d/N$, $c_L=N$, and all other $c_j=1$.
Take a first row in the direction $A_t/\sqrt{Q_t}$ and choose positive
factor norms
\[
 s_j=r/\sqrt{c_j},\qquad
 r=(dP^2Q_t)^{1/(2L)}.
\]
Their product is $P\sqrt{Q_t}$, and each energy term is $r^2$.
The constructed weak-margin parameter therefore has energy
\begin{equation}\label{eq:Ut}
 U_t=L\left(\frac{dQ_t}{4\alpha^{2m}}\right)^{1/L}.
\end{equation}
Slightly increasing the readout gives an open strict-margin region of
consistent wrong predictors, with energies arbitrarily close to $U_t$.

We next show that every consistent weak-margin predictor with the
correct, negative sign at $q_{t+1}$ has energy strictly greater than
$U_t$. Suppose instead that one has energy at most $U_t$.
Lemma~\ref{lem:linearize} bounds its linearization error by
\begin{equation}\label{eq:delta}
 \delta=\frac{\epsilon m\sqrt{pQ_t}}{2\alpha^m}
             \le\frac1{128pQ_*}.
\end{equation}
Indeed, Bernoulli's inequality gives
$\alpha^m\ge1-m\epsilon\ge127/128>1/2$, and
$Q_t\le Q_*$. Put $r_0=1/2-\delta>0$.
The affine vector $B=A_\theta/r_0$ satisfies every observed unit-margin
constraint and $B(q_{t+1})\le-1$. By \eqref{eq:affine-gap} and
\eqref{eq:linear-energy},
\begin{equation}\label{eq:correct-lower}
 \kappa_{d,N}(\theta)\ge
    L\bigl(d r_0^2(Q_t+4/p)\bigr)^{1/L}.
\end{equation}
This exceeds \eqref{eq:Ut}. To check the strict inequality, set
$h=4/(pQ_t)$; here $0<h\le1/2$. Then
\[
 (1-2\delta)^2(1+h)
 \ge1+h-6\delta
 \ge1+\frac4{pQ_*}-\frac6{128pQ_*}
 >1+\frac1{16pQ_*}.
\]
On the other hand, $\epsilon\le1/2$ and $m\epsilon\le1/128$ give
\[
 \alpha^{-2m}\le\exp(4m\epsilon)\le1+8m\epsilon
             \le1+\frac1{16pQ_*}.
\]
The last inequality uses $pQ_*\ge1$ in \eqref{eq:epsilon}.
Thus $(1-2\delta)^2(Q_t+4/p)>Q_t\alpha^{-2m}$, contradicting
\eqref{eq:correct-lower} and $\kappa_{d,N}\le U_t$.

The weak-margin set fitting the history and correctly labeling the
next query is closed. At each fixed width its quadratic energy is
coercive, so its infimum is attained if the set is nonempty.
The preceding strict exclusion therefore gives an energy gap above
$U_t$. Lemma~\ref{lem:gaussian} shows that the correct-label posterior
probability tends to zero. The finite-round argument following that
lemma gives all three prediction conclusions.

For the teacher, at depth $k+2$ use $A_M/2^k$ as the first affine vector,
then $k$ scalar hidden weights equal to two, and readout one.
Every first coefficient is at most one in magnitude. Add hidden scalar
weights one if $L>k+2$. The output has the sign of $A_M$, which is $f$,
and magnitude at least $\alpha^{L-1}|A_M(\xi)|>1/2$.
Every coefficient is at most two. Zero padding is valid because
$\phi(0)=0$, hence $C_L(f)\le2$.

If $C_L(f)=1$, all layers would be scalar and all coefficients at most
one. The downstream scalar map would be $1$-Lipschitz and vanish at zero.
Its sign, up to one overall reversal, would be the first affine sign.
A strict output margin would force a first affine margin $\rho>1/2$.
Normalize that vector by $\rho$ and use the carry inequalities in
Lemma~\ref{lem:selector}. Its original top-bit antisymmetric coefficient
would exceed $2^{k-2}>1$, contradicting the coefficient bound one.
Thus $C_L(f)=2$.
\end{proof}

\begin{corollary}[A larger perturbation with exponentially many errors]
\label{cor:near-half}
Retain the architecture, target, and fixed input permutation of
Theorem~\ref{thm:near}, but count only the rounds
$q_{2^{k-1}},\ldots,q_M$. All three posterior mistake conclusions hold
on these $2^{k-1}$ rounds under the weaker sufficient condition
\begin{equation}\label{eq:epsilon-half}
 0<\epsilon\le
 \frac{2^{k+1}}{128(L-1)\sqrt{d+1}\,Q_*^{3/2}}.
\end{equation}
The teacher still has $C_L(f)=2$ for every $L\ge k+2$, and the
activation condition remains independent of learner width. The upper
bound on $\epsilon$ is of order $4^{-k}/[(L-1)\sqrt{k}]$.
\end{corollary}
\begin{proof}
Put $t_0=2^{k-1}-1$ and $a=4(t_0+1)=2^{k+1}$. For a history through
$q_t$, with $t_0\le t<M$, adding the correct negative label at
$q_{t+1}$ gives the next selector history. Its minimum affine
squared norm is exactly $Q_{t+1}$ by Lemma~\ref{lem:selector}.
Consequently the additional squared-norm cost is
\[
 \Delta_t=Q_{t+1}-Q_t=4(t+1)\ge a.
\]
This improves the coarse bound $4/p$ used in the theorem's proof.

The same scalar candidate has energy $U_t$ from \eqref{eq:Ut}.
Since $a\le Q_*$, condition \eqref{eq:epsilon-half} gives
$m\epsilon\le1/128$ and $\alpha^m>1/2$ as before. A parameter of
energy at most $U_t$ has linearization error at most
\begin{equation}\label{eq:delta-half}
 \delta=\frac{\epsilon m\sqrt{pQ_t}}{2\alpha^m}
       \le\frac{a}{128Q_*}.
\end{equation}
If it correctly labels the next query, dividing its affine
linearization by $1/2-\delta$ gives a vector of squared norm at least
$Q_{t+1}=Q_t+\Delta_t$. Set $h=\Delta_t/Q_t$.
Here $0<h<1/2$: indeed $k\ge3$, $t\ge t_0\ge3$, and
\[
 Q_t-8(t+1)\ge35+2t(t-3)>0.
\]
Thus
\[
 (1-2\delta)^2(1+h)
 \ge1+h-6\delta
 \ge1+\frac{a}{Q_*}-\frac{6a}{128Q_*}
 >1+\frac{a}{2Q_*},
\]
whereas
\[
 \alpha^{-2m}\le1+8m\epsilon
 \le1+\frac{a}{16\sqrt p\,Q_*^{3/2}}
 \le1+\frac{a}{16Q_*}.
\]
The linearization energy bound therefore excludes every correct
next-query parameter from the energy ball at $U_t$. The same closed-set
energy gap and Gaussian argument prove the posterior conclusions.
Finally $\alpha^{L-1}>1/2$ gives the identical bounded teacher, and
the exclusion of $C_L=1$ is unchanged.
\end{proof}
